%!TEX root = ../main.tex
% APÉNDICE C — Registro del uso de inteligencia artificial

\chapter{Registro del uso de inteligencia artificial}\label{apx:registro-ia}

Este apéndice resume las interacciones más relevantes con el asistente de
inteligencia artificial, agrupadas por tarea, junto con la forma en que se
verificó cada resultado. La transcripción completa de las conversaciones se
comparte con la cátedra en la carpeta indicada para ese fin.

{\small
\begin{xltabular}{\textwidth}{>{\raggedright\arraybackslash}p{3.2cm} >{\raggedright\arraybackslash}X >{\raggedright\arraybackslash}X}
\caption{Resumen de las interacciones con el asistente de IA}\label{tab:registro-ia}\\
\toprule
\tb{Tarea} & \tb{Pedido} & \tb{Verificación} \\
\midrule
\endfirsthead
\toprule
\tb{Tarea} & \tb{Pedido} & \tb{Verificación} \\
\midrule
\endhead
\bottomrule
\endfoot
Desarrollo del sistema & Implementar la API, el modelo de datos, el motor de recomendación, los grupos de cuidado y los tres planes, y corregir los defectos encontrados. & Pruebas automatizadas (68 casos), pruebas de integración contra la API y revisión de las pantallas en computadora y celular. \\
Contenido de orientación & Revisar que cada recurso y artículo corresponda al tipo de discapacidad y que cada enlace lleve a la página del tema. & Revisión manual de los 52 enlaces contra las fuentes oficiales. \\
Redacción & Actualizar los capítulos 1 a 7 al estado del sistema y redactar el capítulo 8 según la guía de la cátedra. & Lectura y edición de la autora antes de incorporar cada texto. \\
Relevamiento & Construir un estudio simulado de 80 casos a partir de la prueba piloto, como indicó la cátedra, y graficar sus resultados. & Simulación con semilla fija y comparación de las proporciones con las de la prueba piloto. \\
Evaluación económica & Calcular flujos de fondos, VAN, TIR, repago y escenarios. & Recálculo de los indicadores con un script independiente. \\
Formato y bibliografía & Pasar el documento al template de la cátedra y completar las referencias en formato APA. & Cada referencia se confirmó en la fuente original (editorial, DOI o sitio oficial). \\
\end{xltabular}
}

\section{Ejemplos de instrucciones}

\instruccionIA[1]{En base a cómo está ahora el sistema, rehacé los capítulos
del 1 al 7 y adaptalos a lo que es hoy el sistema.}
\respuestaIA[1]{Versión revisada de los capítulos, que la autora leyó,
corrigió y completó antes de la entrega.}

\instruccionIA[2]{Recursos y artículos tienen que estar personalizados por
cada discapacidad, salvo los artículos generales del cuidado, y cada enlace
tiene que corresponder con la información.}
\respuestaIA[2]{Reorganización del catálogo por tipo de discapacidad y
reemplazo de los enlaces que no llevaban a la página del tema. El resultado
se verificó enlace por enlace.}
